\chapter{Conclusiones}
\label{cap:conclusiones}

Esta tesis se propuso mejorar la escalabilidad de la síntesis composicional de controladores
atacando su cuello de botella: el paso de minimización que se ejecuta una vez por iteración y
que hasta ahora se resolvía calculando \textit{Weak Synthesis Observational Equivalence} (WSOE)
con un algoritmo de complejidad $\bigo{nm}$. La propuesta fue reemplazarlo por
\textit{branching bisimilarity}, para la cual existe el algoritmo $\bigo{m \log n}$ de Groote,
Jansen, Keiren y Wijs~\cite{groote2019}.

El primer resultado del trabajo fue descubrir que el reemplazo no es directo. Las dos
equivalencias no coinciden, y lo que hace falta garantizar es una dirección concreta: que
\textit{branching} nunca fusione un par de estados que WSOE mantiene separado, ya que eso
significaría reducir de más y potencialmente perder soluciones del problema de control. La
hipótesis natural ---tomar como acciones internas a todas las acciones locales--- resulta
incorrecta, y el Ejemplo~3 de la Sección~\ref{sec:ejemplos} exhibe un contraejemplo concreto: al
ocultar una acción controlable local, \textit{branching} fusiona un par que WSOE distingue.
Restringiendo las acciones internas a las acciones locales \textit{no controlables}
($\tau = \Upsilon_u$) la correspondencia se restablece, y el
Teorema~\ref{thm:wsoe-refina-branching} demuestra que en ese caso vale
$\sim_{B} \ \C \ \sim_{W}$. La inclusión además es estricta, como muestra el Ejemplo~4 adaptado
de van Glabbeek y Weijland~\cite{glabbeek1996}.

\vspace{.5cm}
Con ese resultado como base, los objetivos alcanzados fueron los siguientes:

\begin{itemize}
    \item Caracterizar formalmente la relación entre WSOE y \textit{branching bisimilarity} en
    el contexto de la síntesis composicional, identificando la condición sobre el conjunto de
    acciones internas bajo la cual el reemplazo es correcto.

    \item Implementar el algoritmo $\bigo{m \log n}$ de Groote et al.\ dentro de \MTSA{}
    \cite{mtsa} e integrarlo al flujo composicional en reemplazo de la llamada existente a
    WSOE. La implementación atravesó tres iteraciones: una primera versión funcional, una
    segunda con optimizaciones parciales, y una tercera que incorpora las estructuras
    \texttt{RefinablePartition} y \texttt{LinkedTransitionPartitions} necesarias para alcanzar
    efectivamente la cota teórica.

    \item Desarrollar una metodología de evaluación basada en la captura de $2212$ instancias
    reales de minimización, extraídas instrumentando una corrida completa de la síntesis
    composicional sobre el conjunto de casos de prueba existente. Esto permitió iterar sobre la
    implementación sin pagar el costo de reejecutar la síntesis en cada experimento.

    \item Validar la correctitud de la implementación contra dos oráculos independientes: la
    implementación de referencia de mCRL2~\cite{mcrl2}, escrita por los propios autores del
    algoritmo, y el resultado del algoritmo de WSOE, verificando sobre cada instancia que nunca
    ocurriera la fusión que el Teorema~\ref{thm:wsoe-refina-branching} prohíbe.

    \item Evaluar empíricamente el rendimiento. La versión final se mantiene por debajo de los
    $1000$~ms en las instancias más grandes, donde WSOE alcanza los $7600$~ms, y la ventaja se
    acentúa sobre los modelos con mayor cantidad de transiciones internas. En el $83\%$ de las
    instancias ambas equivalencias producen exactamente la misma reducción.
\end{itemize}

\vspace{.5cm}
El trabajo también deja identificadas algunas limitaciones:

\begin{itemize}
    \item \textbf{Uso de memoria.} Contra lo que sugiere el análisis del paper original, todas
    las versiones de \textit{branching} consumieron en promedio más memoria que WSOE. La razón
    es estructural: WSOE mantiene una única partición de estados y recalcula el resto en cada
    paso, mientras que el algoritmo de Groote et al.\ mantiene simultáneamente particiones de
    estados y de transiciones con sus índices asociados. Es un intercambio explícito de memoria
    por tiempo.

    \item \textbf{La mejora sólo se materializa en instancias grandes.} Por debajo de los
    $n \approx 100$ estados el tiempo está dominado por costos fijos y no hay diferencia
    apreciable entre los algoritmos. Como la mayoría de las instancias que genera la síntesis
    composicional son chicas, el impacto sobre el tiempo total del método depende de cuántas
    instancias grandes aparezcan en cada problema.

    \item \textbf{Las optimizaciones no son opcionales.} Las versiones v1 y v2, que implementan
    el mismo algoritmo pero sin las estructuras de datos que el paper especifica, no superan a
    WSOE. Recién la v3 materializa la mejora asintótica. Esto confirma que la cota
    $\bigo{m \log n}$ es una propiedad del algoritmo \textit{junto con} sus estructuras, y no
    del esquema de refinamiento en abstracto.
\end{itemize}

\vspace{.5cm}
Como trabajo a futuro, presentamos las siguientes ideas:

\begin{itemize}

    \item Reducir el consumo de memoria de la implementación, explorando representaciones más
    compactas de las particiones de transiciones o el uso de arreglos primitivos en lugar de las
    estructuras de colecciones de Java. (aca poner el nuevo paper)

    \item Rep. simbolica
\end{itemize}
